% Separate title page for the Digital Creativity submission (double-anonymous review)
\documentclass[12pt]{article}
\usepackage[a4paper,margin=2.7cm]{geometry}
\usepackage{fontspec}
\setmainfont{Times New Roman}
\pagestyle{empty}
\begin{document}
\begin{center}
{\LARGE\bfseries The critic and the builder:\\ how the Chuzhditsa papers were written}\\[2em]
{\large Anastas Stoyanovsky}\\[0.5em]
Unaffiliated (independent researcher)\\
\texttt{anastas@stoyanovs.ky}
\end{center}
\vspace{2em}
\noindent\textbf{Article type:} Research article. \quad \textbf{Word count:} approx.\ 9{,}400 words (body, excluding references and the corpus appendix).
\medskip

\noindent\textbf{Declaration of AI use} (per Taylor \& Francis policy): the manuscript was drafted by a large language model (Claude, Anthropic) at the direction and under the editorial control of the human author, who is solely accountable for its content. This use is not incidental: the AI-mediated production process is the object of study, and Section ``Authorship and disclosure'' of the manuscript states the position in full, including the operational disclosure (tool, dates, manner of use, human verification). No AI tool is listed as an author.
\medskip

\noindent\textbf{Data availability:} the critique corpus, audit script, and complete artifact history are public in the project repository (reference withheld from the anonymized manuscript).
\medskip

\noindent\textbf{Author biography} (under 100 words): Anastas Stoyanovsky is a mathematician and engineer. Trained in pure mathematics (MSc, Purdue University), he has spent a decade building natural-language and information-retrieval systems — at IBM Watson and Amazon among others — and holds patents in natural-language processing. As an independent researcher he designed and directed the Chuzhditsa project: a phonological citation register for Cyrillic and the four-family parametric typeface that implements it.
\end{document}
